\documentclass{article}
\usepackage{tcee-repaso}
\tcetema[La OMC: servicios, propiedad intelectual y plurilaterales]{3.B.34}{La OMC. Los acuerdos distintos de los de mercancías.}
\tcefecha{14/11/2026}
\begin{document}
\cabecerarepaso
\begin{columnas}

% ==========================================================================
\bloque{Esquema del cante}
\esquemacante{%
\apartadocante{Introducción: «ojos que no ven, corazón que no siente»}{2,5}
\apartadocante{I. El Acuerdo General sobre el Comercio de Servicios}{10}
\subcante{1. Delimitación e importancia}{2}
\subcante{2. Motivación económica: barreras regulatorias}{2,5}
\subcante{3. Estructura del AGCS}{3}
\subcante{4. Evolución de las negociaciones}{1,5}
\subcante{5. Valoración y perspectivas}{1}
\apartadocante{II. El Acuerdo sobre los ADPIC}{9,5}
\subcante{1. Delimitación e importancia}{1,5}
\subcante{2. Motivación económica: Norte y Sur}{2,5}
\subcante{3. Estructura del ADPIC}{2,5}
\subcante{4. De la OMPI a la salud pública}{2}
\subcante{5. Valoración y perspectivas}{1}
\apartadocante{III. Acuerdos plurilaterales y otros acuerdos}{7}
\subcante{1. Contratación pública}{2,5}
\subcante{2. Otros acuerdos e iniciativas}{2,5}
\subcante{3. Transparencia y diferencias}{1}
\subcante{4. Situación actual}{1}
\apartadocante{Conclusión}{1}}
\emph{Gráficos del cante}: el mapa de los acuerdos, los modos de suministro y los umbrales modo~1/modo~3 (\hyperref[rep:I]{I}), el de la protección Norte-Sur (\hyperref[rep:II]{II}) y el de la preferencia nacional en las compras públicas (\hyperref[rep:III]{III}); fórmulas: el orden de los umbrales y $\bar T^{\,*}=v(L_N+L_S)/(b+2d)$.

% ==========================================================================
\bloque[2,5]{Introducción}
\begin{itemize}
\item \textbf{Enganche}: «Ojos que no ven, corazón que no siente»: los servicios eran las \emph{partidas invisibles}, y sus barreras son regulaciones, no aranceles; las ideas, buena parte del valor de un móvil o un medicamento. Yaundé (marzo de 2026): expira por primera vez desde 1998 la moratoria de aranceles a las transmisiones electrónicas. ¿Se pueden poner reglas comunes a un comercio que no se ve?
\item \textbf{Relevancia}: servicios 9,6 billones de \$ (2025), 27,5\,\% del comercio de bienes y servicios, casi la mitad en valor añadido; cobros por propiedad intelectual $>$600.000 M\,\$. España: servicios, ingeniería y construcción, indicaciones geográficas; negocia la UE \vertema{3.B.41}.
\item \textbf{Contextualización}: OMC (1-1-1995) \vertema{3.B.33}. Distintos de mercancías: anexo 1B \concepto{AGCS}, 1C \concepto{ADPIC}, 2 ESD (lo propio), 3 \concepto{MEPC}, 4 \concepto{plurilaterales} (ACP) e iniciativas conjuntas. Ronda Uruguay: servicios y patentes a cambio de agricultura y textiles.
\item \textbf{Problemática}: ¿por qué un acuerdo si las barreras no son aranceles? ¿Debe un acuerdo comercial obligar a proteger, y a quién beneficia? ¿Cómo avanzar sin consenso?
\item \textbf{Estructura}: \hyperref[rep:I]{I} AGCS; \hyperref[rep:II]{II} ADPIC; \hyperref[rep:III]{III} plurilaterales y otros. Mismo patrón: delimitación, motivación, estructura, negociaciones, valoración.
\end{itemize}

\graficorepaso[\linewidth]{acuerdos-mapa}{El mapa: anexos 1B, 1C, 2, 3 y 4, y las iniciativas conjuntas fuera de los anexos}

% ==========================================================================
\bloque[10]{I. El AGCS\label{rep:I}}
\begin{ideaclave}
Reglas multilaterales para barreras que son regulaciones: 4 modos, obligaciones generales y compromisos en listas positivas. Gran mérito, la transparencia; límite, liberalizó poco más de lo que ya estaba abierto.
\end{ideaclave}

\sub{I.1. Delimitación e importancia\label{rep:I.1}}
\begin{itemize}
\item \concepto{Servicio}: cambia el estado de quien lo recibe; no se almacena; producción y consumo a la vez $\Rightarrow$ alguien tiene que moverse.
\item \concepto{Modos} (art.~1.2): 1 transfronterizo (35\,\% en 2022), 2 consumo en el extranjero (8\,\%), 3 presencia comercial (56\,\%, inversión \vertema{3.B.18}), 4 personas físicas (1\,\%). El modo~1 crece: 26\,\% en 2005.
\item \emph{Cuantitativa}: 2/3 del PIB y del empleo; 27,5\,\% de las exportaciones (21\,\% en 2005); +8\,\% en 2025; prestados digitalmente 5,3 billones (55\,\%). \concepto{Servificación}: casi el 50\,\% en valor añadido (TiVA); \concepto{curva de la sonrisa} (\autor{Stan Shih}).
\item \emph{Cualitativa}: primeras reglas; insumos de todo; intereses ofensivos (PD) y defensivos (PED, salvo modo~4).
\end{itemize}

\graficorepaso[\linewidth]{modos-suministro}{Quién cruza la frontera en cada modo y su peso (TiSMoS, 2022)}

\graficorepaso[\linewidth]{servicios-sector}{Servicios comerciales por sector: de 2,6 a 9,6 billones de \$ (2005-2025)}

\sub{I.2. Motivación: barreras regulatorias como costes fijos\label{rep:I.2}}
\begin{itemize}
\item \emph{Idea}: ventaja comparativa, escala y \autor{Melitz} valen para los servicios; la barrera es regulación (discriminatoria o heterogénea) que se paga una vez por mercado: coste fijo. Modo~1 = exportar; modo~3 = filial (\autor{Helpman}, \autor{Melitz} y \autor{Yeaple}, 2004; \autor{Kox} y \autor{Nordås}, 2007).
\item \emph{Supuestos}: \autor{Melitz}; modo~1 con $F^{\mathrm{M1}}=F^x+R^{\mathrm{M1}}$ e iceberg $\tau$; modo~3 con $F^{\mathrm{M3}}=F^{\mathrm I}+R^{\mathrm{M3}}$, sin $\tau$.
\item \emph{Desarrollo}: $B^{\mathrm{M1}}=\tau^{1-\sigma}\Omega\varphi^{\sigma-1}-wF^{\mathrm{M1}}$; $B^{\mathrm{M3}}=\Omega\varphi^{\sigma-1}-wF^{\mathrm{M3}}$.
\begin{formula}
\funcionclave{F^{\mathrm{M3}}>\tau^{\sigma-1}F^{\mathrm{M1}}>F^d\Rightarrow\varphi^*<\varphi^*_{\mathrm{M1}}<\varphi^*_{\mathrm{M3}}}
\end{formula}
Bajar $R^{\mathrm{M3}}$ (acceso en modo~3) baja $\varphi^*_{\mathrm{M3}}$.
\item \emph{Implicaciones}: margen extensivo (pocas empresas, grandes); importa la heterogeneidad, no solo la discriminación (reconocimiento); el modo~3 domina por la proximidad; sin ingresos arancelarios.
\item \emph{Extensiones}: costes de comerciar servicios $\approx$2 veces los de bienes (OMC, 2019); modo~4 $\approx$ migración. \emph{Evidencia}: STRI (aéreo 0,40, jurídicos 0,35, contabilidad 0,33; España, entre las más abiertas); \autor{Nordås} y \autor{Rouzet} (2017); \autor{Breinlich} y \autor{Criscuolo} (2011); \autor{Arnold}, \autor{Javorcik} y \autor{Mattoo} (2011). \emph{Valoración}: la regulación tiene fines legítimos $\Rightarrow$ regular bien, no desregular.
\end{itemize}

\graficorepaso[0.85\linewidth]{servicios-modos-hmy}{Cierran, solo interior, modo~1, modo~3; liberalizar el modo~3 baja su umbral}

\graficorepaso[0.9\linewidth]{stri-sectores}{STRI 2025: media, recorrido y España por sector}

\sub{I.3. Estructura del AGCS\label{rep:I.3}}
\begin{itemize}
\item 29 arts., 6 partes. \emph{Alcance}: todo salvo facultades gubernamentales y derechos de tráfico aéreo; cubrir $\neq$ liberalizar (excepción cultural). Clasificación W/120: 12 sectores, $\approx$160 subsectores.
\item \emph{Obligaciones generales}: \concepto{NMF} (art.~II, con exenciones); \concepto{transparencia} (art.~III); \concepto{reglamentación nacional} (art.~VI); \concepto{reconocimiento} (art.~VII); monopolios (art.~VIII); pagos (art.~XI). \emph{Excepciones}: integración (art.~V), balanza de pagos (art.~XII), generales (art.~XIV), seguridad (art.~XIV bis), cautelares financieras. Sin disciplinas: salvaguardias (art.~X), contratación (art.~XIII), subvenciones (art.~XV).
\item \emph{Compromisos específicos}: \concepto{acceso a los mercados} (XVI, 6 limitaciones: proveedores, valor, operaciones, personas, persona jurídica, capital extranjero) y \concepto{trato nacional} (XVII, \emph{negociable}). \concepto{Lista positiva}; «ninguna» / «sin consolidar»; horizontales.
\item \emph{Anexos}: financieros (\emph{prudential carve-out}; 5.º Protocolo, 1999, 70 miembros); telecomunicaciones (acceso a redes; 4.º Protocolo, 1998, 69 gobiernos; Documento de Referencia) \vertema{4.A.12}; personas físicas; transporte aéreo y marítimo.
\end{itemize}

\sub{I.4. Negociaciones\label{rep:I.4}}
\begin{itemize}
\item Ronda Uruguay: consolidar lo abierto; 1997: telecomunicaciones y finanzas. Doha (desde 2000): peticiones y ofertas; bloqueada con la agricultura.
\item Exención PMA (2011; hasta 2030). \emph{TiSA}: 23 partes (50 países), suspendido en 2016. \concepto{Reglamentación nacional} (dic. 2021, 67 miembros, $>$90\,\% del comercio; por las listas y la NMF; hoy 73 participantes, en vigor desde 2024).
\end{itemize}

\sub{I.5. Valoración\label{rep:I.5}}
\begin{itemize}
\item \emph{A favor}: alcance, flexibilidad, transparencia, juez (\emph{Juegos de azar}, DS285). \emph{En contra}: consolida más que abre; polizón por la NMF; incompleto; atado a la agricultura. Futuro: plurilaterales y servicios digitales.
\end{itemize}

% ==========================================================================
\bloque[9,5]{II. El Acuerdo sobre los ADPIC\label{rep:II}}
\begin{ideaclave}
No obliga a abrir, sino a \emph{proteger}, con mínimos y juez. Sin acuerdo se protege demasiado poco; con mínimos iguales, el Sur paga: de ahí las flexibilidades de la salud pública.
\end{ideaclave}

\sub{II.1. Delimitación e importancia\label{rep:II.1}}
\begin{itemize}
\item Derechos de autor y conexos; propiedad industrial. Objetivo (art.~7): innovación y difusión, equilibrio productores-usuarios.
\item UE: sectores intensivos en PI 47\,\% del PIB (2017-2019); cobros mundiales 607.000 M\,\$ (2025): EE.~UU. +134; Irlanda $-$191; falsificaciones 2,3\,\% del comercio (2021).
\end{itemize}

\graficorepaso[\linewidth]{propiedad-intelectual}{Cobros, pagos y saldo por propiedad intelectual (2025)}

\sub{II.2. Motivación: la protección entre el Norte y el Sur\label{rep:II.2}}
\begin{itemize}
\item \emph{Idea}: ideas no rivales (\autor{Romer}, 1990) \vertema{3.A.44}; patente óptima (\autor{Nordhaus}, 1969); protección del país = bien público mundial $\Rightarrow$ polizón (\autor{Grossman} y \autor{Lai}, 2004).
\item \emph{Supuestos}: $N$ inventa; $T_i$ protección; $\bar T=L_NT_N+L_ST_S$; $n=\kappa b\bar T$; $W_S=L_Sn(v-(b+d)T_S)$; $W_N$ + beneficio de sus inventores.
\item \emph{Desarrollo}: $T_S=\frac{L_Sv-(b+d)L_NT_N}{2(b+d)L_S}$ (sustitutivas); $T_N=\frac{L_Nv-dL_ST_S}{(b+2d)L_N}$. Ejemplo: Nash $T_N=0{,}48$, $T_S=0{,}10$.
\begin{formula}
\funcionclave{\bar T^{\,*}=\frac{v(L_N+L_S)}{b+2d}}
\end{formula}
Óptimo 1 frente a 0,57: poca protección. Armonizar ($H$): mundo 0,82 $\to$ 1; Sur 0,49 $\to$ 0,25; Norte 0,33 $\to$ 0,75.
\item \emph{Implicaciones}: motivo para el acuerdo; el Norte protege más; la eficiencia \emph{no} exige armonizar; paquete con agricultura y textiles; flexibilidades.
\item \emph{Extensiones}: \autor{Helpman} (1993): el Sur pierde, el Norte puede perder; \autor{Lai} (1998): por IED sí más innovación; \autor{Scotchmer} (1991): innovación acumulativa. \emph{Evidencia}: \autor{McCalman} (2001); \autor{Branstetter}, \autor{Fisman} y \autor{Foley} (2006); \autor{Kyle} y \autor{McGahan} (2012). \emph{Valoración}: un solo número; importan las excepciones.
\end{itemize}

\graficorepaso[0.72\linewidth]{grossman-lai}{Nash $E$ por debajo de la recta eficiente; $H$ (armonizar) o $A$ (protege el Norte)}

\sub{II.3. Estructura del ADPIC\label{rep:II.3}}
\begin{itemize}
\item 73 arts., 7 partes. \emph{Principios}: mínimos (París 1883, Berna 1886); \concepto{trato nacional} y \concepto{NMF}; \concepto{protección equilibrada} (7-8); \concepto{agotamiento} libre (art.~6: importaciones paralelas).
\item \emph{Plazos}: patentes 20 años; autor 50; marcas 7 renovables; dibujos 10; esquemas de trazado 10; IG y secretos sin plazo.
\item \emph{Patentes}: excluibles orden público, métodos médicos, plantas y animales (no microorganismos); \concepto{excepción Bolar} (art.~30; DS114); \concepto{licencias obligatorias} (art.~31: negociar antes, remuneración, \emph{principalmente} mercado interno).
\item \emph{IG}: protección reforzada de vinos y bebidas (art.~23: ni «tipo Rioja»).
\item \emph{Observancia}: civil, frontera, penal. \emph{Transición}: 1996 (PD), 2000/2005 (PED), PMA hasta 1-7-2034 (farmacéuticos 1-1-2033).
\end{itemize}

\sub{II.4. De la OMPI a la salud pública\label{rep:II.4}}
\begin{itemize}
\item OMPI (1967/1970; ONU 1974): sin juez. EE.~UU. (sección~301 especial) lleva la PI a Punta del Este.
\item \concepto{Declaración de Doha} (14-11-2001): apoyar la salud pública; libertad para licencias obligatorias y agotamiento. Párrafo~6: decisión de 2003 $\to$ \concepto{art.~31 bis} (2005; en vigor 23-1-2017; 142 aceptaciones); Canadá-Ruanda (2008).
\item COVID-19: India y Sudáfrica (oct. 2020); EE.~UU. solo vacunas (mayo 2021); la UE, licencias obligatorias. \concepto{CM12} (17-6-2022): vacunas, exportar sin art.~31.f, 5 años; extensión a diagnósticos y tratamientos: no acordada (2022-2024); nunca se ha usado.
\item Abiertas: registro de IG; divulgación del origen (tratado OMPI 2024); reclamaciones sin infracción (moratoria expirada en 2026); ADPIC-plus; ACTA (rechazado por el PE, 4-7-2012).
\end{itemize}

\sub{II.5. Valoración\label{rep:II.5}}
\begin{itemize}
\item \emph{A favor}: unifica, da juez, flexibilidad (Doha, 31 bis, CM12). \emph{En contra}: mínimos que paga el Sur; estático (datos, IA); observancia desigual (China: DS362, DS611).
\end{itemize}

% ==========================================================================
\bloque[7]{III. Plurilaterales y otros acuerdos\label{rep:III}}
\begin{ideaclave}
Geometría variable: lo nuevo avanza entre los que quieren. Límites: el polizón si se extiende por la NMF y el veto si pide entrar en el anexo~4.
\end{ideaclave}

\sub{III.1. Contratación pública\label{rep:III.1}}
\begin{itemize}
\item Fuera del GATT (III.8) y del AGCS (XIII). \concepto{ACP}: Tokio (1979), 1994 (en vigor 1996), \emph{revisado} (30-3-2012; en vigor 6-4-2014). 22 partes, 49 miembros; Australia (2019), Reino Unido (2021), Macedonia del Norte (2023). Compras 12-13\,\% del PIB (OCDE); $>$1,7 billones de \$ abiertos al año.
\item \emph{Sesgo nacional} (\autor{Baldwin}, 1970): país pequeño, $p^*$, el gobierno compra $G$ solo a nacionales.
\begin{formula}
\funcionclave{G\leq S^d(p^*)\Rightarrow\Delta p=\Delta M=0;\ G>S^d(p^*)\Rightarrow p_g>p^*}
\end{formula}
Muerde donde el gobierno es casi el único comprador (obras, defensa). \emph{Extensiones}: \autor{McAfee} y \autor{McMillan} (1989); \autor{Trionfetti} (2000). \emph{Evidencia}: sesgo grande en la UE (\autor{Herz} y \autor{Varela-Irimia}, 2020).
\end{itemize}

\graficorepaso[0.75\linewidth]{contratacion-sesgo}{Neutral con $G_1$; con $G_2$, $p_g>p^*$ y sobrecoste}

\sub{III.2. Otros acuerdos e iniciativas\label{rep:III.2}}
\begin{itemize}
\item \concepto{MIC}: contenido local y equilibrio prohibidos (arts.~III y XI); 2/5/7 años; Canadá (2013), India (2016), DS623.
\item \concepto{IFD}: 2017-2023; procedimientos, no protección; 131 participantes; sin consenso en Abu Dabi y Yaundé (165 de 166 a favor).
\item \concepto{ATI}: 1996, 29 $\to$ 84 participantes, 97\,\%; 2015: 201 productos, 24 participantes (53 miembros).
\item Moratoria digital (1998-2026). \concepto{Comercio Electrónico}: texto de julio de 2024 (sin flujos de datos); texto cerrado en dic. 2024; Yaundé: 66 (hoy 68) lo aplicarán con 45 aceptaciones \vertema{4.A.14}.
\item Pymes (\emph{Global Trade Helpdesk}); medio ambiente (Bienes Ambientales suspendido; TESSD, plásticos, combustibles fósiles); competencia fuera desde 2004.
\end{itemize}

\sub{III.3. Transparencia y diferencias\label{rep:III.3}}
\begin{itemize}
\item \concepto{MEPC}: cada 3/5/7 años; último de la UE, abril de 2023; el siguiente, diciembre de 2026 exámenes. ESD \vertema{3.B.33}: \concepto{represalias cruzadas} en el ADPIC (Ecuador 2000, Antigua 2007, Brasil 2009); no infracción; MPIA (DS611).
\end{itemize}

\sub{III.4. Situación actual\label{rep:III.4}}
\begin{itemize}
\item Multilaterales congelados; Yaundé sin resultados en servicios y PI. Actividad en plurilaterales; reforma (toma de decisiones). La UE, en todas las iniciativas, ACP y MPIA \vertema{3.B.41}.
\end{itemize}

% ==========================================================================
\bloque[1]{Conclusión}
\begin{itemize}
\item \textbf{Recapitulación}: 2 comercios invisibles: servicios (regulaciones como costes fijos) e ideas (protección con reparto desigual); lo nuevo, plurilateral.
\item \textbf{Relevancia}: para España importan tanto como los aranceles.
\item \textbf{Extensiones}: \vertema{3.B.33} \vertema{3.B.35} \vertema{3.B.41} \vertema{3.B.18} \vertema{3.A.44} \vertema{4.A.7} \vertema{4.A.22}.
\item \textbf{Opinión}: valen más por la transparencia que por lo que liberalizaron; plurilaterales abiertos y sin veto.
\item \textbf{Idea final}: las barreras invisibles eran cómodas; el logro fue ponerlas a la vista. Se pueden poner reglas a lo que no se ve si antes se hace visible.
\end{itemize}

% ==========================================================================
\bloque{Autores y fechas}
\begin{itemize}
\item 1883 París; 1886 Berna; 1967 OMPI; 1969 \autor{Nordhaus}; 1970 \autor{Baldwin}; 1979 código de contratación; 1989 \autor{McAfee}-\autor{McMillan}; 1990 \autor{Romer}; 1993 \autor{Helpman}; 1994 AGCS y ADPIC; 1996 ATI; 1997 protocolos 4.º y 5.º; 2001 Doha y salud pública; 2004 \autor{Grossman}-\autor{Lai} y \autor{Helpman}-\autor{Melitz}-\autor{Yeaple}; 2012 ACP revisado; 2015 ATI ampliado; 2017 art.~31 bis en vigor; 2021 reglamentación nacional; 2022 CM12 (vacunas); 2023 IFD; 2024 Comercio Electrónico; 2026 Yaundé.
\end{itemize}

\bloque{Datos clave}
\begin{itemize}
\item Servicios 2025: 9,6 billones de \$ (+8\,\%), 27,5\,\% del total; digitales 5,3 (55\,\%). Modos (2022): 35/8/56/1\,\%.
\item STRI 2025: aéreo 0,40; jurídicos 0,35; contabilidad 0,33; distribución 0,20.
\item PI 2025: cobros 607.000 M\,\$; EE.~UU. +134; Irlanda $-$191; UE 47\,\% del PIB.
\item ACP: 22 partes, 49 miembros. ATI: 84. IFD: 131. Reglamentación nacional: 73 participantes, en vigor desde 2024.
\end{itemize}

\bloque{Preguntas del tribunal}
\begin{enumerate}
\item \textbf{¿Por qué 4 modos?} Porque el servicio exige proximidad: se mueve el servicio, el cliente, la empresa o la persona.
\item \textbf{¿En qué se diferencia el trato nacional del AGCS del del GATT?} En el AGCS se negocia sector a sector; en el GATT es general.
\item \textbf{¿Lista positiva o negativa?} Positiva: solo lo listado (AGCS); negativa: todo salvo reservas (CETA).
\item \textbf{¿Por qué protege más el Norte?} Se queda con el beneficio de sus inventores: solo paga la pérdida irrecuperable.
\item \textbf{¿Hace falta armonizar para la eficiencia?} No: basta el mercado protegido total (\autor{Grossman} y \autor{Lai}).
\item \textbf{¿Qué es el art.~31 bis?} Licencia obligatoria para exportar a países sin capacidad; primera enmienda de la OMC (2017).
\item \textbf{¿Qué son las represalias cruzadas?} Suspender obligaciones de otro acuerdo (ADPIC): Antigua contra EE.~UU.
\item \textbf{¿Distorsiona la preferencia nacional en las compras?} No si el gobierno compra menos de lo que produce el país (\autor{Baldwin}).
\item \textbf{¿Por qué no entra el IFD en el anexo~4?} El art.~X.9 exige consenso: un solo miembro basta para impedirlo (la India en Yaundé, según la prensa).
\end{enumerate}

\begin{trampa}
\begin{itemize}
\item AGCS en el anexo \emph{1B}; ADPIC en el \emph{1C}; MIC en el \emph{1A} (mercancías).
\item El ATI \emph{no} es del anexo~4: está en las listas y se extiende por la NMF.
\item Derecho de autor: \emph{50} años; patentes: \emph{20} desde la solicitud; marcas: \emph{7} renovables.
\item Art.~31 bis: aprobado en \emph{2005}, en vigor en \emph{2017}.
\item La exención de la CM12 cubre solo \emph{vacunas}.
\item El modo~4 \emph{no} cubre a quien busca empleo ni la residencia.
\end{itemize}
\end{trampa}

\end{columnas}
\end{document}
